\documentclass[UTF8,openany,zihao=5,fontset=none]{ctexbook}

\usepackage[
  a4paper,
  left=2.45cm,
  right=2.45cm,
  top=2.35cm,
  bottom=2.55cm,
  headheight=14pt
]{geometry}
\usepackage{amsmath}
\usepackage{amssymb}
\usepackage{circuitikz}
\usepackage{enumitem}
\usepackage{fancyhdr}
\usepackage{fontspec}
\usepackage{graphicx}
\usepackage[obeyspaces,spaces]{url}
\usepackage{hyperref}
\usepackage{listings}
\usepackage{longtable}
\usepackage{microtype}
\usepackage{pgfplots}
\usepackage{rotating}
\usepackage{tabularx}
\usepackage[most]{tcolorbox}
\usepackage{tikz}
\usepackage{xcolor}
\usepackage{xeCJK}
\usepackage{siunitx}

\usetikzlibrary{
  arrows.meta,
  calc,
  circuits.logic.US,
  decorations.pathreplacing,
  fit,
  matrix,
  positioning,
  shapes.geometric
}
\pgfplotsset{compat=1.18}
\sisetup{
  detect-all,
  group-separator={\,},
  group-minimum-digits=4,
  per-mode=symbol
}

\definecolor{BookInk}{HTML}{17202A}
\definecolor{BookMuted}{HTML}{5D6D7E}
\definecolor{BookNavy}{HTML}{17365D}
\definecolor{BookBlue}{HTML}{2471A3}
\definecolor{BookTeal}{HTML}{148F77}
\definecolor{BookAmber}{HTML}{B9770E}
\definecolor{BookRed}{HTML}{B03A2E}
\definecolor{BookCodeBg}{HTML}{F4F6F7}
\definecolor{BookRule}{HTML}{CCD1D1}

\defaultfontfeatures{Ligatures=TeX}
% 正文、标题和代码承担不同职责。正文使用衬线体保证长篇阅读节奏，标题使用
% 无衬线体建立层级，Maple Mono 只用于代码、寄存器、路径和日志。
\IfFontExistsTF{TeX Gyre Pagella}{
  \setmainfont{TeX Gyre Pagella}
}{
  \setmainfont{Latin Modern Roman}
}
\IfFontExistsTF{TeX Gyre Heros}{
  \setsansfont{TeX Gyre Heros}
}{
  \setsansfont{Latin Modern Sans}
}
\setCJKmainfont{FandolSong}[AutoFakeBold=2,AutoFakeSlant=0.2]
\setCJKsansfont{FandolHei}[AutoFakeBold=2]
\IfFontExistsTF{Maple Mono NF CN}{
  \setmonofont{Maple Mono NF CN}[Ligatures=NoCommon]
  \setCJKmonofont{Maple Mono NF CN}[AutoFakeBold=2,FallBack=FandolFang]
}{
  \IfFontExistsTF{Maple Mono NL NF CN}{
    \setmonofont{Maple Mono NL NF CN}[Ligatures=NoCommon]
    \setCJKmonofont{Maple Mono NL NF CN}[AutoFakeBold=2,FallBack=FandolFang]
  }{
    \setmonofont{Latin Modern Mono}[Ligatures=NoCommon]
    \setCJKmonofont{FandolFang}[AutoFakeBold=2]
  }
}

\linespread{1.08}
\setlength{\parindent}{2em}
\setlength{\parskip}{0.08em}
\setlength{\tabcolsep}{4pt}
\raggedbottom
\sloppy
\emergencystretch=2em

% 三位页码需要比 LaTeX 默认值更宽的目录页码盒。
\makeatletter
\renewcommand{\@pnumwidth}{2.1em}
\makeatother

\ctexset{
  chapter={
    format=\Huge\bfseries\sffamily\color{BookNavy},
    titleformat=\Huge\bfseries\sffamily\color{BookNavy}
  },
  section={format=\Large\bfseries\sffamily\color{BookBlue}},
  subsection={format=\large\bfseries\sffamily\color{BookNavy}}
}

\pagestyle{fancy}
\fancyhf{}
\renewcommand{\chaptermark}[1]{\markboth{\thechapter\quad #1}{}}
\renewcommand{\sectionmark}[1]{\markright{\thesection\quad #1}}
\fancyhead[LE]{\footnotesize\sffamily\color{BookMuted}\thepage\quad\nouppercase{\leftmark}}
\fancyhead[RO]{\footnotesize\sffamily\color{BookMuted}\nouppercase{\rightmark}\quad\thepage}
\renewcommand{\headrulewidth}{0.25pt}
\fancypagestyle{plain}{
  \fancyhf{}
  \fancyfoot[C]{\small\color{BookMuted}\thepage}
  \renewcommand{\headrulewidth}{0pt}
}

\hypersetup{
  bookmarksdepth=2,
  colorlinks=true,
  linkcolor=BookBlue,
  urlcolor=BookBlue,
  citecolor=BookBlue,
  pdftitle={从复位向量到自研 x86-64 操作系统},
  pdfauthor={x86-64 OS Lab},
  pdfsubject={结合完整源码学习 x86-64 操作系统}
}

\DeclareRobustCommand{\code}[1]{\textcolor{BookBlue}{\texttt{\nolinkurl{#1}}}}
\DeclareRobustCommand{\filepath}[1]{\textcolor{BookTeal}{\texttt{\nolinkurl{#1}}}}
\DeclareRobustCommand{\term}[2]{\textbf{#1}（#2）}
\newcommand{\topic}[1]{%
  \par\medskip{\normalsize\bfseries\color{BookNavy}#1}\quad\ignorespaces
}
\newcolumntype{L}[1]{>{\raggedright\arraybackslash}p{#1}}

\lstdefinestyle{oscode}{
  basicstyle=\ttfamily\small,
  backgroundcolor=\color{BookCodeBg},
  frame=single,
  rulecolor=\color{BookRule},
  columns=fullflexible,
  keepspaces=true,
  breaklines=true,
  showstringspaces=false,
  xleftmargin=0.8em,
  xrightmargin=0.4em
}
\lstset{style=oscode}

\newtcolorbox{keyidea}{
  enhanced,
  breakable,
  sharp corners,
  boxrule=0pt,
  leftrule=0.7pt,
  colback=BookCodeBg,
  colframe=BookBlue,
  title=核心思想,
  fonttitle=\bfseries,
  coltitle=BookNavy
}

\newtcolorbox{contractbox}{
  enhanced,
  breakable,
  sharp corners,
  boxrule=0pt,
  leftrule=0.7pt,
  colback=white,
  colframe=BookTeal,
  title=边界契约,
  fonttitle=\bfseries,
  coltitle=BookTeal
}

\newtcolorbox{failurebox}{
  enhanced,
  breakable,
  sharp corners,
  boxrule=0pt,
  leftrule=0.7pt,
  colback=white,
  colframe=BookRed,
  title=失败路径,
  fonttitle=\bfseries,
  coltitle=BookRed
}

\newenvironment{partbridge}
{\clearpage\thispagestyle{plain}\vspace*{0.5em}
 \noindent{\Large\bfseries\color{BookNavy}本部分导读}\par\medskip}
{\par\clearpage}

\newenvironment{pdfdiagram}[1][]
{\par\medskip\noindent\begin{center}\begin{tikzpicture}[os diagram,#1]}
{\end{tikzpicture}\end{center}\medskip}

\newcommand{\diagramnote}[1]{%
  \par{\noindent\footnotesize\color{BookMuted}图示：#1\par}\smallskip
}

\tikzset{
  os diagram/.style={
    font=\footnotesize\sffamily,
    >=Latex,
    node distance=10mm and 14mm
  },
  os block/.style={
    draw=BookBlue,
    fill=BookCodeBg,
    rounded corners=2pt,
    align=center,
    minimum height=9mm,
    inner xsep=6pt
  },
  os hardware/.style={os block,draw=BookAmber,fill=white},
  os state/.style={os block,draw=BookTeal},
  os arrow/.style={-{Stealth[length=2.4mm]},draw=BookMuted,thick}
}


\setcounter{tocdepth}{1}
\setcounter{secnumdepth}{2}

\begin{document}

\frontmatter
\begin{titlepage}
\thispagestyle{empty}
\vspace*{2.2cm}
{\Huge\bfseries\color{BookNavy} 从复位向量到自研\\x86-64 操作系统\par}
\vspace{0.8cm}
{\Large\color{BookBlue} 硬件契约、启动链、内核机制与工程证据\par}
\vspace{1.2cm}
{\large x86-64 OS Lab\par}
\vfill
{\small\color{BookMuted}
从 CPU 上电后的第一条指令开始，不借用 BIOS、UEFI 或现成引导器；
QEMU 只模拟硬件，其余软件逐层自行实现。\par}
\end{titlepage}

\chapter*{前言}

写一个操作系统很容易被误解为“先写一个内核主函数”。真正的机器没有主函数，
也不会替开发者建立栈、清空未初始化区或选择执行模式。CPU 复位以后只承诺一组
架构状态；固件必须从固定地址取到合法指令，初始化最早的观测通道，找到后续
程序，把处理器带入长模式，最后才可能进入 C++20 内核。

本书和配套工程坚持一个严格边界：QEMU 只模拟 x86-64 CPU、内存、串口、IDE
等硬件，不使用 SeaBIOS、OVMF、GRUB、Limine、Multiboot 或
\code{-kernel} 替代启动链。汇编统一采用 NASM Intel 语法，高层代码采用
freestanding C++20。构建和测试脚本运行在宿主机，可以使用 Python；进入
镜像的目标代码则必须由我们掌握。

全书既讲原理，也讲证据。每个阶段都回答四个问题：输入状态是什么，输出状态
是什么，失败怎样暴露，验收怎样自动重复。读者不必在开篇掌握全部硬件知识，
但应当愿意追踪地址、寄存器、字节布局和控制流，而不是依赖“能跑就算完成”。

\tableofcontents

\mainmatter
\part{方法、边界与硬件契约}
\begin{partbridge}
这一部分先确定“自己写操作系统”的准确边界，再建立 CPU 复位、地址、总线和
QEMU 硬件模型。它回答代码执行之前必须成立的条件。
\end{partbridge}
\chapter{把操作系统学习变成可验收工程}

\section{学习对象不是代码量}

操作系统连接处理器、内存和设备，同时向程序提供受保护的抽象。学习项目真正
要获得的是机制之间的因果关系：复位向量为什么必须位于固定地址；页表为什么
同时承担翻译和权限；中断为什么改变执行流；调度器为什么不能只维护一个函数
指针数组。代码只是把这些关系变成可执行证据。

\begin{keyidea}
每次只增加一个新的硬件或软件不变量。先让 ROM 可执行，再让串口可观察，再
加载下一阶段；不要同时引入模式切换、ELF、页表和 C++ 运行时。
\end{keyidea}

\section{项目边界}

本项目允许 QEMU TCG 模拟硬件，因而宿主机不需要是 x86-64。QEMU 不提供
启动软件；传给 \code{-bios} 的 128 KiB 镜像完全由项目生成。宿主 Python
负责调度 CMake、CTest、NASM、LLD 和 QEMU，但不会进入目标镜像。

\begin{center}
\begin{tikzpicture}[os diagram]
  \node[os block] (source) {NASM / C++20\\项目源代码};
  \node[os state,right=of source] (image) {ROM / 磁盘\\可审计镜像};
  \node[os hardware,right=of image] (qemu) {QEMU TCG\\硬件模型};
  \node[os state,below=of qemu] (evidence) {串口、状态、\\测试证据};
  \draw[os arrow] (source) -- (image);
  \draw[os arrow] (image) -- (qemu);
  \draw[os arrow] (qemu) -- (evidence);
\end{tikzpicture}
\end{center}

\section{软件工程闭环}

需求定义“必须做到什么”；架构说明模块和交接；ADR 记录关键取舍；代码实现
机制；测试保存可重复证据；发布记录对应一个可复现里程碑。完成阶段必须同时
满足构建、静态审计、宿主测试、整机测试、失败路径和文档六个方面。

\section{语言与代码规则}

目标代码使用显式宽度类型，如 \code{std::uint32_t} 和
\code{std::uint64_t}。硬件寄存器按协议宽度表示；内部计数、容量和地址若
没有更窄的硬件约束，优先使用 64 位。类成员访问显式写 \code{this->}；
常量采用 \code{OS_模块_功能} 的全大写命名；头文件与实现分离；中文注释
解释约束和原因，不复述语法。

\chapter{x86 复位状态与第一条指令}

\section{复位是契约，不是初始化完成}

上电复位后，处理器处于实地址模式。最关键的可观察状态是指令入口：
\code{CS} 的可见值为 \code{0xF000}，隐藏基址为
\code{0xFFFF0000}，\code{IP/EIP} 指向 \code{0xFFF0}，因此第一条
指令的物理地址是：
\[
  0xFFFF0000 + 0xFFF0 = 0xFFFFFFF0.
\]
隐藏基址是理解现代 x86 复位地址的关键。此时不能简单套用普通实模式
\(segment \times 16 + offset\) 去解释第一次取指。

\section{128 KiB ROM 的顶端映射}

项目 ROM 覆盖 \code{0xFFFE0000} 到 \code{0xFFFFFFFF}。文件偏移与物理
地址之间满足：
\[
  physical = 0xFFFE0000 + file\_offset.
\]
所以复位向量的文件偏移是 \code{0x1FFF0}，入口
\code{0xFFFFF000} 的文件偏移是 \code{0x1F000}。

\begin{center}
\begin{tikzpicture}[os diagram]
  \node[os hardware,minimum width=42mm] (rom) {ROM 起点\\0xFFFE0000};
  \node[os state,above=12mm of rom,minimum width=42mm] (entry)
    {16 位入口\\0xFFFFF000};
  \node[os block,above=12mm of entry,minimum width=42mm] (reset)
    {复位向量\\0xFFFFFFF0};
  \node[os hardware,above=8mm of reset,minimum width=42mm] (end)
    {ROM 末端\\0xFFFFFFFF};
  \draw[os arrow] (rom) -- node[right]{128 KiB 地址空间} (entry);
  \draw[os arrow] (entry) -- (reset);
  \draw[os arrow] (reset) -- (end);
\end{tikzpicture}
\end{center}

\section{为什么使用 16 位近跳转}

复位向量放置操作码 \code{E9} 和有符号 16 位位移。近跳转只修改 IP，不
重载 CS，因而保留隐藏基址 \code{0xFFFF0000}。目标 IP 是
\code{0xF000}，物理目标正好是 \code{0xFFFFF000}。位移为：
\[
  0xF000 - (0xFFF0 + 3) = -4083 = 0xF00D.
\]
因此复位向量前三个字节是 \code{E9 0D F0}。

\begin{failurebox}
如果使用会重载 CS 的远跳转，就必须重新证明目标段如何映射到 ROM；如果链接
脚本把入口放错一个字节，CPU 也不会给出友好错误，只会执行填充值或触发异常。
\end{failurebox}

\chapter{平台、地址空间与 I/O 端口}

\section{CPU 地址不等于 RAM}

物理地址空间由平台译码。某段地址可以落入 DRAM、ROM、PCI 窗口或空洞。
固件所在的高地址 ROM 和低地址 RAM 同时存在；把栈设为 \code{0x7000}
是平台布局决策，不是 x86 指令自动保证的事实。

\section{端口映射 I/O}

传统 PC 串口使用独立 I/O 端口空间，由 \code{in}、\code{out} 指令访问。
COM1 基址是 \code{0x3F8}。基址加偏移选择寄存器：数据寄存器偏移 0，
中断使能偏移 1，线路控制偏移 3，线路状态偏移 5。某些偏移的含义还会受
DLAB 位影响，因此寄存器访问是一段状态协议，不能看成互不相关的赋值。

\begin{contractbox}
软件写寄存器前必须知道设备状态；设备返回状态位后，软件才能推进所有权。
轮询必须有界，错误路径必须有可观察结果。
\end{contractbox}

\section{QEMU 在本项目中的职责}

QEMU 提供 \code{pc} 机器、\code{qemu64} CPU、TCG 执行、RAM、ISA 串口和
IDE 磁盘模型。它不替固件初始化串口，不替引导器读磁盘，也不替内核建立页表。
\code{-nodefaults} 减少无关默认设备，\code{-display none} 和
\code{-serial stdio} 把最早日志绑定到测试进程。


\part{从源文件到第一条串口日志}
\begin{partbridge}
这一部分把 NASM、链接脚本、ROM 字节、复位向量和 16550A 串口连成一条可
逐字节审计的启动路径，并说明怎样证明成功与失败。
\end{partbridge}
\chapter{工具链怎样制造可执行事实}

\section{三个不同层次}

NASM 把 Intel 语法汇编成 ELF32 可重定位目标文件；LLD 按链接脚本为各节
分配最终虚拟地址；\code{llvm-objcopy} 去掉 ELF 元数据并用
\code{0xFF} 填充地址空洞，生成 QEMU 读取的裸 ROM。每一步解决不同问题：
指令编码、地址决议、介质字节布局。

\begin{center}
\begin{tikzpicture}[os diagram]
  \node[os block] (asm) {.asm\\符号与指令};
  \node[os state,right=of asm] (obj) {.o\\可重定位节};
  \node[os block,right=of obj] (elf) {.elf\\最终地址};
  \node[os state,right=of elf] (bin) {.bin\\128 KiB 字节};
  \draw[os arrow] (asm) -- node[above]{NASM} (obj);
  \draw[os arrow] (obj) -- node[above]{LLD} (elf);
  \draw[os arrow] (elf) -- node[above]{objcopy} (bin);
\end{tikzpicture}
\end{center}

\section{链接脚本是硬件布局的一部分}

普通应用程序由操作系统装载器解释 ELF；ROM 没有装载器。链接脚本必须直接
表达 ROM 起点、入口、复位向量和末端，并用 \code{ASSERT} 阻止入口代码覆盖
最后 16 字节。地址选择一旦错误，源码语义再正确也无法执行。

\section{为什么固件先使用 ELF32}

当前入口只包含 16 位指令，但 NASM 的 \code{elf32} 输出能让 LLD 处理节和
符号；它不表示 CPU 已进入 32 位模式。执行模式由 CPU 状态和指令前缀决定，
目标文件容器格式与运行模式必须分开理解。

\section{宿主与目标}

工具可能运行在 AArch64、x86-64 或其他宿主架构，产物仍由 NASM 和 LLD
明确指向 x86。QEMU TCG 翻译目标指令，因此构建机和模拟 CPU 不必同构。

\chapter{ROM 布局与复位代码走读}

\section{入口先建立确定状态}

入口依次执行 \code{cli}、\code{cld}，清零数据段、附加段和栈段，设置栈顶，
再初始化串口。关闭中断是因为 IDT 尚不存在；清除方向标志保证字符串读取向
高地址推进；显式设置栈是因为 \code{call}、\code{ret} 和
\code{push} 已经依赖它。

\section{为什么字符串使用 CS 覆盖}

固件代码和常量位于高地址 ROM。入口的 CS 隐藏基址仍指向 ROM 高区，而 DS
被清零并用于低地址 RAM。使用 \code{cs lodsb} 可以按当前代码段读取日志
字符串，避免假设 ROM 在低端还存在别名。

\section{源码走读}

下面的清单直接引用配套工程，书与代码保持同源：

\lstinputlisting[
  language={[x86masm]Assembler},
  firstline=1,
  lastline=177,
  caption={16 位复位与串口入口}
]{../../../../source/firmware/src/reset_and_serial.asm}

\section{ROM 审计}

宿主工具不通过“文件存在”判断成功。它检查镜像恰为 131072 字节，复位偏移的
操作码为 \code{E9}，解码位移后目标为 \code{0xF000}，入口以
\code{CLI/CLD} 开始。这是对最终字节的独立验证，不依赖汇编源码意图。

\chapter{16550A 串口：第一个观测通道}

\section{初始化顺序}

固件先关闭 UART 中断，设置 DLAB，写入波特率除数 1，再关闭 DLAB 并配置
8 数据位、无校验、1 停止位。随后配置 FIFO 和 modem control。对 QEMU
模拟串口而言，宿主终端不真正按 115200 波特限制吞吐，但寄存器协议仍应按照
真实设备编写。

\begin{longtable}{p{0.24\textwidth}p{0.20\textwidth}p{0.46\textwidth}}
\toprule
动作 & 端口/值 & 含义 \\
\midrule
关闭中断 & 0x3F9 / 0x00 & 最早阶段采用轮询，不接收 IRQ \\
打开 DLAB & 0x3FB / 0x80 & 让偏移 0、1 指向除数锁存器 \\
设置除数 & 0x3F8 / 1，0x3F9 / 0 & 115200 基准波特率 \\
设置 8N1 & 0x3FB / 0x03 & 同时关闭 DLAB \\
设置 FIFO & 0x3FA / 0xC7 & 启用并清空 FIFO \\
设置 modem & 0x3FC / 0x0B & 拉起 DTR、RTS、OUT2 \\
\bottomrule
\end{longtable}

\section{发送不是直接写字节}

写数据前读取线路状态寄存器，等待 bit 5 表示发送保持寄存器为空。若直接连续
写而不检查，代码隐含了设备永远足够快的假设。当前实现用 16 位计数器进行最多
\code{0xFFFF} 次轮询，成功通过进位标志返回，超时清除进位标志。

\section{两条日志的协议意义}

\code{RESET} 证明复位跳转、栈和 UART 初始化已经工作；
\code{SERIAL_READY} 证明完整字符串发送路径仍然可用。失败镜像在第一条消息
之后人为屏蔽就绪状态，用同一轮询函数走超时路径。因此测试既能识别“完全未
执行”，也能识别“执行了但设备没有完成”。

\chapter{测试、故障注入与可观察性}

\section{测试金字塔必须落到硬件边界}

宿主单元测试验证地址范围、ROM 解码和串口协议解析；集成测试验证链接产物；
固定种子随机测试生成大量错误复位位移，确认审计器不会误接受；QEMU 系统测试
执行真实目标字节。它们不是重复劳动，而是分别缩小不同故障的定位范围。

\section{有界运行}

固件最终执行 \code{hlt}，符合当前没有后续任务的状态。测试进程给 QEMU
两秒预算，超时后读取捕获的串口内容。成功用例要求两个标记都存在；失败用例
要求 \code{RESET} 存在且 \code{SERIAL_READY} 不存在。异常提前退出同样
视为失败。

\begin{failurebox}
仅检查进程“还活着”不能证明 CPU 执行了我们的代码；仅检查一个字符串也不能
区分初始化成功与发送路径中途失效。可观察协议应让不同执行阶段留下不同证据。
\end{failurebox}

\section{随机测试的边界}

随机测试不替代规格样例。固定种子保证失败可重现；生成器只改变复位位移并保持
其他布局合法；断言是“任何不指向 0xF000 的位移都必须被拒绝”。这种测试
检验的是性质，而不是堆积随机输入。

\section{完整验证命令}

\begin{lstlisting}[language=bash]
python3 tools/os.py verify
python3 tools/os.py test --layer unit
python3 tools/os.py test --layer integration
python3 tools/os.py test --layer randomized
python3 tools/os.py test --layer system
python3 tools/os.py test --layer failure-path
\end{lstlisting}


\part{从实模式到内核}
\begin{partbridge}
这一部分面向后续实现，完整解释保护模式、长模式、ELF64 装载、异常、内存、
中断与设备驱动之间的依赖关系。
\end{partbridge}
\chapter{从实模式走向 64 位内核}

\section{转换不是一条指令}

进入长模式至少需要验证 CPU 能力、建立 GDT、打开 A20、建立页表、设置
CR4.PAE、写 EFER.LME、设置 CR0.PG 和 CR0.PE，再通过远跳转装载 64 位代码
段。顺序错误会导致通用保护异常或三重故障。

\begin{center}
\begin{tikzpicture}[os diagram]
  \node[os state] (real) {16 位实模式};
  \node[os block,right=of real] (protected) {32 位保护模式};
  \node[os block,right=of protected] (paging) {PAE 四级页表};
  \node[os state,right=of paging] (long) {64 位长模式};
  \draw[os arrow] (real) -- (protected);
  \draw[os arrow] (protected) -- (paging);
  \draw[os arrow] (paging) -- (long);
\end{tikzpicture}
\end{center}

\section{A20 与地址环绕}

早期 PC 为兼容 8086 曾让地址第 20 位可屏蔽。若 A20 未开启，1 MiB 以上地址
可能回绕。项目必须自行启用并验证 A20，不能因为 QEMU 默认环境“似乎可用”
就删除这项契约。

\section{Stage 1 与磁盘}

v0.2 将通过 IDE PIO 读取自研 Stage 1。驱动需要选择通道、设备和 LBA，
发送命令，轮询 BSY/DRQ/ERR，并按 16 位数据端口读取扇区。每次传输都要有
超时、扇区范围和错误状态验证。

\section{ELF64 装载}

ELF 头不是值得盲目信任的结构体。装载器验证 magic、类别、端序、机器类型、
程序头表范围和每个 PT\_LOAD 段；检查文件范围不越界、
\code{filesz <= memsz}、目标区间不重叠保留内存；复制文件字节并清零 BSS，
最后构造 BootInfo 交给内核入口。

\chapter{内核入口、异常与内存管理}

\section{内核入口契约}

64 位入口要明确栈对齐、方向标志、中断状态、参数寄存器、页表和可写内存范围。
汇编适配层建立 C++ ABI 后才调用 \code{kernelMain}。freestanding 环境没有
宿主标准库、异常展开和默认运行时，链接审计必须拒绝未知未解析符号。

\section{先处理异常，再增加复杂性}

GDT 定义代码段和数据段，TSS 提供特权切换所需栈，IDT 把向量映射到门描述符。
每个异常桩统一保存寄存器并规范化错误码，再交给 C++ 分发器。页故障日志至少
包含 CR2、错误码和 RIP，否则后续内存管理无法定位。

\section{物理页与虚拟地址空间}

物理页分配器管理“哪些页可用”；页表管理“虚拟地址怎样映射、具有什么权限”；
内核堆管理“小对象怎样复用页面”。三者不能混成一个分配器。

\begin{center}
\begin{tikzpicture}[os diagram]
  \node[os block] (heap) {内核堆\\对象分配};
  \node[os state,below=of heap] (vmm) {虚拟内存\\页表与权限};
  \node[os hardware,below=of vmm] (pmm) {物理页\\所有权位图};
  \draw[os arrow] (heap) -- (vmm);
  \draw[os arrow] (vmm) -- (pmm);
\end{tikzpicture}
\end{center}

\section{大类型与精确宽度}

物理地址、虚拟地址和全局计数优先用 64 位，为后续扩展保留空间；页表项必须
严格用 \code{std::uint64_t}，设备寄存器则按 8/16/32 位协议宽度访问。
“能用大的就用大的”不意味着把端口寄存器错误地扩大，而是避免内部模型被
无意义的小上限锁死。

\chapter{中断、设备与用户边界}

\section{从轮询到中断}

轮询由 CPU 主动反复读取状态；中断由设备在状态变化时请求处理。启用中断前，
内核必须完成 IDT、控制器屏蔽、EOI 规则和共享状态同步。中断上半部只确认和
记录事件，较长工作应延后执行。

\section{PIC、PIT、键盘与 IDE}

PIC 重映射避免 IRQ 与 CPU 异常向量冲突；PIT 提供周期时钟；键盘控制器提供
扫描码；IDE 提供扇区 I/O。每个驱动都应分为寄存器访问、状态机、策略和公共
接口，防止端口魔法数字散落在调度器或文件系统中。

\section{Ring 3 与系统调用}

用户态不能直接访问内核页和特权指令。进入 Ring 3 需要用户代码/数据段、
用户栈、页表权限和正确的返回帧。系统调用入口必须先保存不可信上下文，再验证
编号、指针、长度和权限；验证通过才触碰内核对象。

\begin{contractbox}
用户指针不是 C++ 指针承诺，而是“可能映射、可能跨页、可能在验证后变化”的
不可信地址。复制进出用户空间必须定义部分失败和页故障处理。
\end{contractbox}

\section{驱动的完成语义}

命令写入设备不等于操作完成。驱动必须区分提交、设备接受、数据可用、完成确认
和错误恢复。这个模型会一直延伸到块层、文件系统和用户 \code{read} 返回值。


\part{把最小内核发展成操作系统}
\begin{partbridge}
最后把执行流、同步、文件、系统调用和工程治理组织成长期可演进的学习项目，
并给出从 v0.0 到 v1.0 的完整验收地图。
\end{partbridge}
\chapter{进程、同步与文件系统}

\section{进程是资源所有权集合}

进程包含地址空间、线程、打开文件、权限和生命周期状态。线程保存可调度执行
上下文。调度器在时钟中断或阻塞点选择可运行线程，并保存/恢复寄存器与栈。
抢占发生在任意允许中断的位置，因此共享状态必须有同步规则。

\section{锁、等待与 IPC}

自旋锁适合极短且不能睡眠的临界区；等待队列把不能继续的线程从运行队列移出；
管道把字节缓冲、读写端引用和唤醒条件组合成 IPC 对象。任何同步原语都要说明
谁能睡眠、谁能唤醒、锁顺序和中断上下文约束。

\section{从块到文件}

块设备只认识扇区。文件系统用超级块、inode、目录项和数据块提供名字与层次；
VFS 用统一对象接口连接进程文件描述符和具体文件系统。写操作还需要区分：
复制进缓存、提交给设备、设备报告完成、数据达到持久介质。

\section{失败与恢复}

磁盘可能在任意写入之间掉电。更新多个结构时，需要日志、写时复制或明确的
恢复算法保证旧状态或新状态至少一个完整。测试应注入越界块号、空间耗尽、
中途失败和损坏元数据，而不仅验证创建一个小文件。

\chapter{十三个里程碑与实践方法}

\section{完整路线}

\begin{longtable}{p{0.11\textwidth}p{0.20\textwidth}p{0.60\textwidth}}
\toprule
版本 & 阶段 & 核心验收 \\
\midrule
v0.0 & 工程基线 & 可复现构建、测试分层、CI 与文档结构 \\
v0.1 & 复位与串口 & 自研 ROM 从复位向量输出串口协议 \\
v0.2 & 固件加载 & IDE PIO 加载自研 Stage 1，并覆盖错误状态 \\
v0.3 & Long Mode & 独立完成 16→32→64 位转换 \\
v0.4 & 内核加载 & 校验并加载 ELF64，交接 BootInfo \\
v0.5 & 内核基础 & GDT、IDT、TSS、异常与 panic \\
v0.6 & 内存管理 & 物理页、虚拟地址空间与内核堆 \\
v0.7 & 中断设备 & PIC、PIT、键盘与 IDE 驱动 \\
v0.8 & 用户边界 & Ring 3、用户 ELF 与系统调用 \\
v0.9 & 进程调度 & 抢占调度与进程生命周期 \\
v0.10 & 同步 IPC & 锁、等待队列、睡眠唤醒与管道 \\
v0.11 & 文件系统 & inode、目录、文件和持久化 \\
v1.0 & 用户环境 & Shell、基础命令和整机回归 \\
\bottomrule
\end{longtable}

\section{每个阶段的固定问题}

\begin{enumerate}[label=\arabic*.]
  \item 硬件或上阶段提供什么输入状态？
  \item 本阶段新建立哪些不变量？
  \item 输出交给谁，地址和 ABI 是什么？
  \item 正常路径怎样观察？
  \item 超时、损坏和非法输入怎样失败？
  \item 哪些单元、集成、随机和整机测试保存证据？
\end{enumerate}

\section{代码走读方法}

先从链接脚本和入口找最终地址，再沿调用链跟踪寄存器和内存所有权；遇到端口
访问就回到设备寄存器协议；遇到跨层调用就写下输入、输出与失败状态；最后对照
测试，确认每个关键分支都能被外部证据区分。

\section{学习记录}

每完成一个机制，写一页实验记录：假设、最小实现、观察结果、失败案例、根因和
后续限制。不要只保存成功截图。真正提高定位能力的是“为什么没有输出”“哪一
个字节错了”“哪个状态位没有到达”。


\backmatter
\chapter*{能力与验收清单}
\addcontentsline{toc}{chapter}{能力与验收清单}

\section*{v0.1 完成后应能解释}

\begin{itemize}
  \item CPU 为什么从物理地址 \code{0xFFFFFFF0} 取第一条指令；
  \item 128 KiB ROM 的文件偏移怎样对应高端物理地址；
  \item 近跳转为什么能保留复位时的 CS 隐藏基址；
  \item NASM、LLD、objcopy 分别解决什么问题；
  \item COM1 的 DLAB、除数、8N1、FIFO 和线路状态位；
  \item 为什么串口轮询必须有界；
  \item ROM 审计、单元、随机和 QEMU 测试各自捕获什么错误。
\end{itemize}

\section*{进入下一阶段前}

\begin{itemize}
  \item 完整执行 \code{python3 tools/os.py verify}；
  \item 能从镜像十六进制数据手算复位目标；
  \item 能在 QEMU 串口日志中区分复位成功和串口超时；
  \item 能说明 v0.2 IDE PIO 的输入、输出、超时和越界条件；
  \item 新设计先写 ADR 和验收条件，再扩展实现。
\end{itemize}

\chapter*{规范与延伸阅读}
\addcontentsline{toc}{chapter}{规范与延伸阅读}

\begin{enumerate}
  \item Intel, \emph{Intel 64 and IA-32 Architectures Software Developer's
  Manual}, Volume 3A，处理器初始化、控制寄存器、保护模式和分页。
  \item Intel, \emph{Intel 64 and IA-32 Architectures Software Developer's
  Manual}, Volume 2，指令编码和语义。
  \item QEMU Project, \emph{System Emulation User's Guide}，机器、固件、
  TCG、字符设备和块设备参数。
  \item IEEE/The Open Group, \emph{ELF gABI}，ELF 文件头、程序头和装载段。
  \item National Semiconductor, \emph{PC16550D Universal
  Asynchronous Receiver/Transmitter With FIFOs}，串口寄存器与状态位。
  \item 配套仓库的 \filepath{docs/adr/}、\filepath{docs/architecture.md}
  和 \filepath{docs/releases/}，用于追踪实现决策与阶段证据。
\end{enumerate}


\end{document}
